\chapter{结论展望}

\section{工作总结}

本文围绕服务器日志异常检测问题，设计并实现了一种基于日志分析和 Liquid 架构的异常检测系统。系统以原始服务器日志为输入，经过日志读取、预处理、模板解析、特征构建、模型训练、阈值校准和结果输出等步骤，实现了窗口级异常检测。

本文主要完成了以下工作。

\begin{enumerate}
    \item \textbf{完成了系统需求分析与总体设计。} 本文分析了服务器日志异常检测的应用场景、功能需求和非功能需求，设计了由数据层、处理层、模型层和输出层组成的系统架构。
    
    \item \textbf{实现了日志数据读取与预处理模块。} 系统能够读取原始日志，抽取时间戳、节点、日志级别、组件、正文和标签等字段，并完成缺失值处理、重复记录过滤和文本规范化。
    
    \item \textbf{实现了日志解析与事件 ID 生成方法。} 系统通过规则和正则表达式识别动态参数，将语义相同但参数不同的日志归并为统一模板，并生成事件 ID 序列。
    
    \item \textbf{构建了面向异常检测的联合特征。} 本文以时间窗口为样本单位，融合事件频次、事件序列和相邻事件时间间隔，使模型能够同时利用局部分布、上下文顺序和非规则时间信息。
    
    \item \textbf{实现了基于 Liquid 架构的异常检测模型。} 模型以联合特征序列为输入，输出窗口级异常评分，并通过验证集阈值校准生成最终异常标签。
    
    \item \textbf{完成了系统测试与实验分析。} 本文从功能测试、训练过程、异常评分、混淆矩阵、模型对比和参数影响等角度验证了系统的可行性。
\end{enumerate}

\section{成果分析}

本文研究形成了以下成果。

\begin{enumerate}
    \item \textbf{形成了完整的日志异常检测流程。} 系统覆盖从原始日志输入到异常结果输出的主要环节，各模块之间的数据接口较为明确。
    
    \item \textbf{提高了日志数据的结构化表达能力。} 通过动态参数替换和模板映射，系统将非结构化日志文本转换为稳定的事件序列，降低了文本稀疏性。
    
    \item \textbf{融合了统计特征、序列特征和时间特征。} 联合特征表示使模型既能感知局部事件分布，又能利用日志事件顺序和真实时间间隔。
    
    \item \textbf{验证了 Liquid 架构在该任务中的可行性。} 实验结果表明，引入时间间隔信息的 Liquid 模型能够完成窗口级异常检测，并在本文实验设置下取得较为稳定的检测效果。
\end{enumerate}

\section{不足分析}

本文系统仍存在以下不足。

\begin{enumerate}
    \item \textbf{日志解析方法依赖人工规则。} 当日志格式变化较大或动态参数类型较复杂时，人工设计的正则规则可能无法覆盖所有情况，模板质量会影响后续检测效果。
    
    \item \textbf{特征语义表达仍然有限。} 本文主要使用模板 ID、频次和时间间隔特征，对日志正文中的深层语义信息利用不足。
    
    \item \textbf{实时处理能力有待提高。} 本文实验以离线数据处理为主，尚未充分验证高并发流式日志场景下的处理延迟和资源消耗。
    
    \item \textbf{泛化能力仍需进一步验证。} 不同系统的日志格式、事件类型和异常模式差异较大，模型迁移到新场景时可能需要重新解析和训练。
    
    \item \textbf{与当前高水平方法仍有差距。} 近年来对比学习、预训练语言模型和大语言模型方法在公开日志数据集上取得了更高结果。本文主要完成本科阶段系统实现和方法验证，尚未充分利用深层语义表示、跨数据集迁移和复杂调参策略。
\end{enumerate}

\section{未来工作}

针对上述不足，后续研究可以从以下方面改进。

\begin{enumerate}
    \item \textbf{改进日志解析方法。} 后续可以引入更自动化的日志模板抽取方法，减少人工规则维护成本，提高对复杂日志格式的适应能力。
    
    \item \textbf{增强日志语义表示。} 可以结合预训练语言模型或领域词向量，将日志模板和正文语义纳入特征表示，以提升模型对语义相近事件的识别能力。
    
    \item \textbf{优化模型结构与训练策略。} 后续可进一步比较 Liquid 架构与其他连续时间模型、注意力机制和图结构模型的结合方式。
    
    \item \textbf{提升系统实时处理能力。} 可以将日志解析、特征构建和异常检测部署为流式处理流程，并对延迟、吞吐量和资源占用进行系统评估。
    
    \item \textbf{扩展多数据集验证。} 后续应在更多类型的服务器日志上进行实验，检验系统在不同日志格式、不同异常比例和不同业务场景下的稳定性。
    
    \item \textbf{采用更严格的实验协议。} 后续可以补充多次随机实验、置信区间、消融实验和跨系统测试，使模型效果评价更加稳健，避免单一划分或单一指标造成结果偏高。
\end{enumerate}

\section{本章小结}

本章对全文工作进行了总结，并分析了本文系统的研究成果、不足和后续改进方向。总体而言，本文围绕服务器日志异常检测构建了较完整的系统流程，并验证了联合特征和 Liquid 架构在日志异常检测任务中的可行性。本文结果应理解为有限实验条件下的系统验证，而不是对当前公开 SOTA 的超越。后续研究仍需在自动化解析、语义建模、实时检测、严格评测和跨场景泛化方面进一步完善。
